% Generated draft: requires mathematical and source review.
\documentclass[11pt]{article}
\usepackage[margin=25mm]{geometry}
\usepackage{fontspec}
\setmainfont{Latin Modern Roman}
\usepackage{amsmath,amssymb,mathtools,unicode-math}
\setmathfont{Latin Modern Math}
\usepackage{enumitem,needspace,xurl,hyperref}
\setlength{\parindent}{0pt}
\setlength{\parskip}{5pt}
\setlength{\emergencystretch}{3em}
\newcommand{\N}{\mathbb{N}}
\newcommand{\Z}{\mathbb{Z}}
\newcommand{\Q}{\mathbb{Q}}
\newcommand{\R}{\mathbb{R}}
\newcommand{\C}{\mathbb{C}}
\newcommand{\F}{\mathbb{F}}
\newcommand{\A}{\mathbb{A}}
\newcommand{\PP}{\mathbb{P}}
\newcommand{\E}{\mathbb{E}}
\newcommand{\eps}{\varepsilon}
\newcommand{\dd}{\,\mathrm{d}}
\begin{document}
\section*{UnsolvedMath ID: 20002689}
\textbf{Source problem code:} catalogue entry
\subsection{Definitions and mathematical statement}
The source takes tensor-independent tuples \(X=(X_1,\ldots,X_k)\), \(Y=(Y_1,\ldots,Y_k)\) whose joint laws are freely infinitely divisible, and asks how scalar Fourier power series from their tensor convolution relate to ordered noncommutative series from free convolution. Freely infinitely divisible means admitting an \(m\)-fold free-convolution root for every integer \(m\ge1\). The record does not state self-adjointness, the state, or a series convention. The aid's precise comparison uses bounded self-adjoint tuples in a noncommutative probability space; tensor independence means the two generated algebras commute and the state factors across them.

For one tuple let \(m_{i_1\cdots i_r}=\varphi(X_{i_1}\cdots X_{i_r})\). Its scalar function has expansion
\[
 \varphi\left(e^{i\sum_jt_jX_j}\right)
 =1+\sum_{r\ge1}\frac{i^r}{r!}
 \sum_{i_1,\ldots,i_r}m_{i_1\cdots i_r}t_{i_1}\cdots t_{i_r}.
\]
The variables \(t_j\) commute, so each coefficient combines all words with the same letter counts. It therefore forgets ordered-word information. For internally noncommuting tuples this scalar function is not automatically the Fourier transform of a classical probability measure on \(\R^k\); the ordinary joint-distribution interpretation requires additional commutativity. Neither case identifies this content-symmetrized function with the full noncommutative \(R\)-series.

Define ordered tensor cumulants \(\kappa\) by the moment expansion over all set partitions, and free cumulants \(r\) by the same expansion over noncrossing partitions, preserving the original order within each block and multiplying scalar block values. Their respective ordered series linearize tensor and free convolution. They agree through degree three. At degree four the sole additional crossing partition is \(\{\{1,3\},\{2,4\}\}\), so for a fixed tuple
\[
 r_{i_1i_2i_3i_4}=\kappa_{i_1i_2i_3i_4}
 +\kappa_{i_1i_3}\kappa_{i_2i_4}.
\]

Couple fixed marginals once tensor-independently and once freely, and let \(A_{ij}=\varphi((X_i-\varphi X_i)(X_j-\varphi X_j))\), and similarly \(B\), be their covariance matrices. The aid's exact fourth ordered moment difference is
\[
 m^{\mathrm{tensor}}_{i_1i_2i_3i_4}
 -m^{\mathrm{free}}_{i_1i_2i_3i_4}
 =A_{i_1i_3}B_{i_2i_4}+B_{i_1i_3}A_{i_2i_4}.
\]
For self-adjoint tuples with a positive state, these covariance matrices are positive semidefinite; the difference vanishes for every length-four word if and only if \(A=0\) or \(B=0\). The comparison does not need infinite divisibility; that is an additional assumption of the original question. It is a quartic diagnostic, not a full classification of obtainable series.

The aid supplies a different known set-level identity due to Belinschi--Nica. Let \(\mathcal D_c(k)\) denote laws of bounded self-adjoint \(k\)-tuples, \(\mathcal D_c^{\mathrm{FID}}(k)\) its freely infinitely divisible subclass, and \(M_\mu\) the ordered moment series without constant term. With Boolean cumulant series
\[
 \eta_\mu=M_\mu(1+M_\mu)^{-1},
\]
the reported Boolean-to-free identity is
\[
 \{R_\mu:\mu\in\mathcal D_c^{\mathrm{FID}}(k)\}
 =\{\eta_\nu:\nu\in\mathcal D_c(k)\}.
\]
This concerns Boolean cumulants, not classical Fourier cumulants. Its full hypotheses and reference are not supplied, and it is not substituted for the source's unresolved request for a precise tensor/Fourier versus free-series comparison.
\subsection{Short English statement}
Scalar Fourier coefficients lose word order, while tensor and free cumulants use different partition families. The aid computes their exact fourth-order convolution defect and separately cites a Boolean-to-free set identity, which is not an identity with classical Fourier series.
\subsection{Sources}
\begin{itemize}
\item Local registry, permanent UnsolvedMath ID \(20002689\), exact public code \texttt{catalogue entry}; canonical location \texttt{aim-probability-notes.json}, \texttt{notes[130]}.
\item AIM, \emph{Free analysis}, subsection 0.10, ``Infinite Divisibility,'' attributed to Nica. The verification field reports prior visual inspection of printed page 7; supplied reference \url{https://aimath.org/WWN/freeanalysis/freeanalysis.pdf}. No PDF was accessed here. No original remarks, year, or literature entry is supplied.
\item Belinschi--Nica Boolean-to-free identity, attributed by the aid without title, date, journal, or theorem number. The ordered quartic comparison and Fourier information-loss statements have no additional primary citation in the record.
\item Supplied research attempt 1: \texttt{partial\_result}, run and registry status \texttt{partially\_solved}, rigor label \texttt{proved}, low novelty confidence; classification \texttt{PARTIAL-PROGRESS}, statement status \texttt{reconstructed\_unverified}, created and updated 14 August 2026. The aid is machine generated, not peer reviewed, and requires independent verification before citation.
\end{itemize}
\end{document}
